% Math/Theory Summary Template (Mode B)
% For pure math explorations without data pipelines.
% Goal: self-contained document that a reader with fresh eyes can follow.
%
% Architecture (house style — see workflow-theory.md "House style" for the whys):
%   - Setup holds PRIMITIVES ONLY (bulleted) + the maintained assumption.
%     Every other object is defined at the top of the section that first uses it.
%   - Each results section: definitions (first use) -> motivation prose ->
%     theorem with FULL inline hypotheses -> strategy paragraph (per-lemma
%     roles) -> assembly proof (forward-citing lemmas) -> shared inputs as
%     numbered displays -> lemmas with bulleted proofs.
%   - One point per line: proofs are itemize/enumerate, one move per bullet.
%   - Key identities as align* chains with && justifications per line.
%   - Explicit summation bounds everywhere; argument-explicit functions in
%     statements; registered notation in statements, raw algebra in proofs.
%   - \clearpage before each section after the first; figures non-floating,
%     never alone on a page.

\documentclass[12pt]{article}

% --- Packages ---
\usepackage{lmodern}
\usepackage[T1]{fontenc}
\usepackage{amsmath,amssymb,amsfonts,mathtools,amsthm}
\usepackage{booktabs}
\usepackage{hyperref}
\usepackage{xcolor}
\usepackage{enumitem}
\setlist{noitemsep, topsep=2pt}
\usepackage{setspace}
\usepackage{tikz}   % figures: draw inline, non-floating
\usepackage{float}

% --- Page Layout ---
\usepackage[margin=1in]{geometry}
\setlength{\parindent}{0pt}
\setlength{\parskip}{0.55em}

% --- Section Formatting ---
\usepackage{titlesec}
\titleformat{\section}{\normalfont\scshape\centering}{\thesection.}{0.5em}{}
\titleformat{\subsection}{\normalfont\bfseries}{\thesubsection.}{0.5em}{}

% --- Hyperlinks ---
\hypersetup{
  colorlinks   = true,
  urlcolor     = blue,
  linkcolor    = blue,
  citecolor    = black,
}

% --- Theorem environments ---
\newtheorem{theorem}{Theorem}
\newtheorem{lemma}[theorem]{Lemma}
\newtheorem{corollary}[theorem]{Corollary}
\newtheorem{definition}[theorem]{Definition}
\theoremstyle{remark}
\newtheorem{remark}[theorem]{Remark}

% --- Project-registered notation ONLY (do not invent shorthand) ---
% \newcommand{...}{...}

\begin{document}

% --- Title Block ---
\begin{center}
{\Large \textbf{PLACEHOLDER: Title}}

{\large PLACEHOLDER: Subtitle}

\medskip

Internal Report \hspace{1em}|\hspace{1em} \today
\end{center}

% --- Plan paragraph: the steps of the document, one sentence each ---
\textbf{Plan.} PLACEHOLDER: one method/idea, then the ordered list of results
(1) ... (2) ... Everything below is proved; the honest boundary is one closing
remark.

%-----------------------------------------------------%
\section{Setup}
%-----------------------------------------------------%

% Primitives ONLY, as bullets. No welfare aggregates, no benchmarks —
% those are defined at first use in their own sections.

PLACEHOLDER: primitives paragraph (type space, payoff, the one derived number
actually used globally).

\begin{itemize}
  \item \textbf{Market.} PLACEHOLDER (weak orderings here; ties dispatched
    inside the proofs that meet them, not assumed away).
  \item \textbf{Equilibrium.} PLACEHOLDER: cite what the source proves
    (sorting, activity, uniqueness) and footnote what is maintained.
\end{itemize}

% Optional inline figure — non-floating, never alone on a page:
% \begin{center}
% \begin{tikzpicture}[x=0.8cm, y=0.8cm, font=\footnotesize] ... \end{tikzpicture}
% \end{center}

%-----------------------------------------------------%
\clearpage
\section{Step 1: PLACEHOLDER first result}
%-----------------------------------------------------%

% 1. Definitions this section needs (first use), argument-explicit:
PLACEHOLDER: The object bounded is ...
\[
\Pi(\pi) \;\equiv\; \dots
\]

% 2. Motivation prose (economics, in words — no bullets here):
PLACEHOLDER: why this result should hold, in words.

% 3. Theorem with FULL inline hypotheses:
\begin{theorem}[PLACEHOLDER]\label{thm:one}
PLACEHOLDER: one-sentence statement; hypotheses inline; conclusion in
registered notation.
\end{theorem}

% 4-5. Proof, bulleted; key identity as an align chain:
\begin{proof}
\begin{itemize}
  \item \emph{PLACEHOLDER first move.}
    \begin{align*}
    X &= \dots && \text{the definition}\\
      &= \dots && \text{insert \eqref{eq:placeholder}}\\
      &= \dots && \text{the key regrouping, named in words}
    \end{align*}
  \item PLACEHOLDER next move, with its reason.
  \item PLACEHOLDER conclusion. \qedhere
\end{itemize}
\end{proof}

%-----------------------------------------------------%
\clearpage
\section{Step 2: PLACEHOLDER main theorem}
%-----------------------------------------------------%

% 1. Definitions at first use (e.g., CS(b_1,...,b_K); benchmarks as values of
%    the same function at special configurations, in words then symbols).

% 2. Motivation prose before the theorem.

% 3. Theorem, full inline hypotheses, complete equality characterization.
\begin{theorem}[PLACEHOLDER]\label{thm:main}
PLACEHOLDER.
\end{theorem}

% 4. Strategy paragraph: name the parts, one clause per lemma's role
%    (bulleted if 3+), flag the crux and where each key assumption enters.

% 5. Assembly proof FIRST (forward-cites the lemmas; exhaustive case list,
%    each case with its own correct reason). May start on a fresh page:
% \clearpage
\begin{proof}[Proof of Theorem~\ref{thm:main}]
\begin{itemize}
  \item PLACEHOLDER case, directly.
  \item PLACEHOLDER case: Lemma~\ref{lem:a}.
  \item PLACEHOLDER case: Lemma~\ref{lem:b}. \qedhere
\end{itemize}
\end{proof}

% 6. Shared inputs as numbered displays (never buried in one lemma's proof):
It remains to prove the lemmas. PLACEHOLDER shared input:
\begin{equation}\label{eq:placeholder}
\dots
\end{equation}

\begin{lemma}[PLACEHOLDER]\label{lem:a}
PLACEHOLDER: full inline hypotheses,
\[
\dots, \qquad \text{for } K \in \{\dots\}.
\]
\end{lemma}

\begin{proof}
\begin{itemize}
  \item PLACEHOLDER, by \eqref{eq:placeholder}.
  \item PLACEHOLDER. \qedhere
\end{itemize}
\end{proof}

\begin{lemma}[PLACEHOLDER]\label{lem:b}
PLACEHOLDER: full inline hypotheses and conclusion.
\end{lemma}

\begin{proof}
\begin{itemize}
  \item PLACEHOLDER, by \eqref{eq:placeholder}.
  \item PLACEHOLDER. \qedhere
\end{itemize}
\end{proof}

%-----------------------------------------------------%
\clearpage
\section*{Closing remark: the honest boundary}
%-----------------------------------------------------%

PLACEHOLDER: what is distribution-free vs. specialized; what remains open
(\textcolor{red}{in red}); the maintained assumption restated in one clause.

\end{document}
